\section{Introducción y objetivos}

El objetivo general del laboratorio es realizar experimentos de clasificación
supervisada con los conjuntos de datos del curso, Potabilidad y Diabetes, usando el
framework de Python construido durante el semestre. Los objetivos específicos son:

\begin{enumerate}
\item Explicar el funcionamiento de KNN, árboles de decisión, Random Forest,
XGBoost y AdaBoost (Sección~2).
\item Definir modelos que combinen preprocesamiento, selección de atributos y cada
algoritmo; para cada uno, comparar la configuración estándar con variaciones,
seleccionar el mejor modelo y analizarlo (Secciones~3 a~5).
\end{enumerate}

Para que las comparaciones sean justas y repetibles se fijaron tres reglas. Primero,
cada variante cambia un solo hiperparámetro respecto del estándar, de modo que
cualquier diferencia se pueda atribuir a ese cambio. Segundo, la elección se hace
con validación cruzada sobre el conjunto de entrenamiento, y el conjunto de prueba
se usa una única vez al final. Tercero, todo el preprocesamiento se aprende sólo
con datos de entrenamiento, para evitar fuga de información.

El documento se organiza así: la Sección~2 presenta los fundamentos teóricos y un
mapa conceptual; la Sección~3, la metodología y su implementación en el framework;
las Secciones~4 y~5, los resultados y el análisis de cada dataset; la Sección~6
compara ambos casos y discute limitaciones, y la Sección~7 resume las conclusiones.
El anexo contiene los comandos para reproducir todo.
